\documentclass[../../../main]{subfiles}

\ifSubfilesClassLoaded{%
    \zexternaldocument*{../../1.Category-Theory/exported/aux/main}%
}{}

\begin{document}
\section{Additive categories}
\begin{definition}[Additive category] \label{def:additive-category}
  A \hyperref[def:category]{category} \(\category{A}\) is \textbf{additive} if it verifies:
  \begin{enumerate}[label=A.\arabic*), ref=A.\arabic*]
    \item\label{axiom:additive-zero} There exists a \hyperref[def:zero-object]{zero object} \(0 \in \category{A}\).
    \item\label{axiom:additive-products} \(\category{A}\) has finite \hyperref[def:categorical-product]{categorical products}.
    \item\label{axiom:additive-hom} For every pair of objects \(X, Y \in \category{A}\) the set \(\Hom_{\category{A}}(X, Y)\) is endowed with an abelian group structure and the composition maps are bilinear with respect to the group operation.
  \end{enumerate}
\end{definition}

\begin{lemma} \label{lemma:products-coproducts-additive}
  Finite \hyperref[def:categorical-product]{products} and \hyperref[def:categorical-coproduct]{coproducts} coincide in an \hyperref[def:additive-category]{additive category}.
\end{lemma}
\begin{proof}
  Consider \(X_{1}, X_{2} \in \category{A}\). We only need to check that \(X_{1} \times X_{2} \cong X_{1} \amalg X_{2}\) since the case for a finite collection of elements follows directly by induction.
  To do this, we will see that \(X_{1} \times X_{2}\) satisfies the universal property of the \hyperref[def:categorical-coproduct]{coproduct}.

  Consider the following diagram
  \begin{center}
    \begin{tikzcd}
      X_{1} \arrow[dd, "\id"'] \arrow[dr, "{(\id, 0)}"] &       & X_{2} \arrow[dd, "\id"] \arrow[dl, "{(0, \id)}"'] \\
        & \arrow[dl, "p_{1}"] X_{1} \times X_{2} \arrow[dr, "p_{2}"'] &   \\
      X_{1} &       & X_{2}
    \end{tikzcd}
  \end{center}

  Note that for any pair of morphisms \((\alpha_{1}, \beta_{1}), (\alpha_{2}, \beta_{2}): R \to X_{1} \times X_{2}\) we have that
  \begin{equation*}
    (\alpha_{1}, \beta_{1}) + (\alpha_{2}, \beta_{2}) = (\alpha_{1} + \alpha_{2}, \beta_{1}, \beta_{2})
  \end{equation*}
  Therefore if \(p_{1}, p_{2}\) are two projections on \(X_{1} \times X_{2}\) whose sum is the identity we have that
  \begin{equation*}
    (\id, 0) \circ p_{1} + (0, \id) \circ p_{2} = \id_{X_{1} \times X_{2}}
  \end{equation*}

  Consider now two morphisms \(f_{i}: X_{i} \to Q\).
  Then its sum \(\phi = f_{1} \circ p_{1} + f_{2} \circ p_{2}\) produces the following diagram
  \begin{center}
    \begin{tikzcd}
      X_{1} \arrow[dd, "\id"'] \arrow[dr, "{(\id, 0)}"] &       & X_{2} \arrow[dd, "\id"] \arrow[dl, "{(0, \id)}"'] \\
        &  X_{1} \times X_{2} \arrow[dd, "\phi"] &   \\
        X_{1} \arrow[dr, "f_{1}"'] &       & \arrow[dl, "f_{2}"] X_{2} \\
              &   Q   &
    \end{tikzcd}
  \end{center}
  Furthermore, this is unique.
  Indeed, consider \(\phi': X_{1} \times X_{2} \to Q\) another morphism filling the diagram. Then
  \begin{align*}
    \phi &= \phi \circ \id_{X_{1} \times X_{2}} \\
      &= \phi \circ ( (\id, 0) \circ p_{1} + (0, \id) \circ p_{2} ) \\
      &= \phi \circ (\id, 0) \circ p_{1} + \phi \circ (0, \id) \circ p_{2} \\
      &= f_{1} \circ p_{1} + f_{2} \circ p_{2} \\
      &= \phi' \circ (\id, 0) + \phi' \circ (0, \id) \\
      &= \phi' \circ ( (\id, 0) \circ p_{1} + (0, \id) \circ p_{2} ) \\
      &= \phi' \circ \id_{X_{1} \times X_{2}} \\
      &= \phi'
  \end{align*}

  Therefore \(X_{1} \times X_{2}\) verifies the universal property of \(X_{1} \amalg X_{2}\).
\end{proof}

\begin{notation}\label{def:direct-sum}
  If \(\{X_{i}\}_{i=1}^{n}\) is a finite set of objects in an \hyperref[def:additive-category]{additive category} we will denote its \hyperref[def:categorical-product]{product} (and therefore by \autoref{lemma:products-coproducts-additive} its \hyperref[def:categorical-coproduct]{coproduct}) by \(\bigoplus_{i=1}^{n} X_{i}\) and call it its \textbf{direct sum}.
\end{notation}

\begin{remark}\label{rem:morphisms-in-direct-sum}
  Let \(X, Y \in \category{A}\) be two objects.
  By the proof of \autoref{lemma:products-coproducts-additive} given morphisms \(p: X \oplus Y \to X, q: X \oplus Y \to Y\) we get morphisms \(i: X \to X \oplus Y, j: Y \to X \oplus Y\) given by the commutative diagram
  \begin{center}
    \begin{tikzcd}
      X \arrow[rr, "\id_{X}"] \arrow[dr, "i"] & & X \\
      & X \oplus Y \arrow[ur, "p"] \arrow[dr, "q"] & \\
      Y \arrow[ur, "j"] \arrow[rr, "\id_{Y}"] & & Y
    \end{tikzcd}
  \end{center}
  Furthermore we have
  \begin{align*}
    p \circ j = q \circ i &= 0 \\
    i \circ p + j \circ q &= \id_{X \oplus Y}
  \end{align*}

  On the other hand, given objects \(X, Y, Z \in \category{A}\) and morphisms \(i, j, p, q\) as above then \(Z = X \oplus Y\).
\end{remark}


\begin{lemma} \label{lemma:mono-epi-additive}
  A morphism \(\varphi: A \to B\) in an \hyperref[def:additive-category]{additive category} is a \hyperref[def:mono-epi-morphism]{monomorphism} if and only if for every morphism \(\xi: Z \to A\)
  \begin{equation*}
    \varphi \circ \xi = 0 \implies \xi = 0
  \end{equation*}
  \hyperref[def:dual]{Dually} \(\varphi\) is an \hyperref[def:mono-epi-morphism]{epimorphism} if and only if for every morphism \(\xi: B \to Z\)
  \begin{equation*}
    \xi \circ \varphi = 0 \implies \xi = 0
  \end{equation*}
\end{lemma}
\begin{proof}
  We will check the \hyperref[def:mono-epi-morphism]{monomorphism} caracterization.
  \begin{itemize}
    \item[\(\implies\)] Let \(\varphi: A \monoarrow B\) be a \hyperref[def:mono-epi-morphism]{monomorphism}.
          Consider a morphism \(\xi: Z \to A\) such that \(\varphi \circ \xi = 0\).
          Since \(\varphi \circ 0 = 0\) applying left-cancellativity we have that \(\xi = 0\).
    \item[\(\impliedby\)] Let \(\xi_{1}, \xi_{2}: Z \to A\) be two morphisms such that \(\varphi \circ \xi_{1} = \varphi \circ \xi_{2}\).

          Since \(\Hom(Z, B)\) is an abelian group we have
          \begin{equation*}
            0 = \varphi \circ \xi_{1} - \varphi \circ \xi_{2} = \varphi \circ (\xi_{1} - \xi_{2})
          \end{equation*}
          Therefore by hypothesis \(\xi_{1} - \xi_{2} = 0 \iff \xi_{1} = \xi_{2}\), proving that \(\varphi\) is a \hyperref[def:mono-epi-morphism]{monomorphism}.
  \end{itemize}
\end{proof}


\begin{definition}[Kernel] \label{def:categorical-kernel}
  Let \(f: A \to B\) be a morphism in an \hyperref[def:additive-category]{additive category}.
  A morphism \(\iota: K \to A\) is a \textbf{kernel} of \(f\) if \(f \circ \iota = 0\) and for every morphism \(g: Z \to A\) such that \(f \circ g = 0\) there exists a unique \(\tilde{g}:Z \to K\) making the following diagram commute
  \begin{center}
    \begin{tikzcd}
      Z \arrow[rr, "g"'] \arrow[dr, dotted, "\tilde{g}"'] \arrow[rrr, bend left, "0"] &   & A \arrow[r, "f"'] & B\\
        & K \arrow[ur, "\iota"'] &   &
    \end{tikzcd}
  \end{center}
\end{definition}

\begin{definition}[Cokernel] \label{def:categorical-cokernel}
  Let \(f: A \to B\) be a morphism in an \hyperref[def:additive-category]{additive category}.
  A morphism \(\pi: B \to C\) is a \textbf{cokernel} of \(f\) if \(\pi \circ f = 0\) and for every morphism \(g: B \to Z\) such that \(\pi \circ g = 0\) there exists a unique \(\tilde{g}:C \to Z\) making the following diagram commute
  \begin{center}
    \begin{tikzcd}
        &   & C \arrow[dr, dotted, "\tilde{g}"] &   \\
      A \arrow[r, "f"] \arrow[rrr, bend right, "0"'] & B \arrow[rr, "g"] \arrow[ur, "\pi"] &   & Z
    \end{tikzcd}
  \end{center}
\end{definition}

\begin{lemma} \label{lemma:kernels-cokernels-are-mono-epi}
  In an \hyperref[def:additive-category]{additive category} \hyperref[def:categorical-kernel]{kernels} are \hyperref[def:mono-epi-morphism]{monomorphisms} and \hyperref[def:categorical-cokernel]{cokernels} are \hyperref[def:mono-epi-morphism]{epimorphisms}.
\end{lemma}
\begin{proof}
  We will only show that \hyperref[def:categorical-cokernel]{cokernels} are \hyperref[def:mono-epi-morphism]{epimorphisms}. The statement about \hyperref[def:categorical-kernel]{kernels} follows analogously.

  Let \(\varphi: A \to B\) be a morphism and let \(\coker \varphi: B \to C\) be its \hyperref[def:categorical-cokernel]{cokernel}.
  Consider another morphism \(\gamma: C \to Z\) such that \(\gamma \circ \coker \varphi = 0\).
  Since \((\gamma \circ \coker \varphi) \circ \varphi = 0\) by definition of the \hyperref[def:categorical-cokernel]{cokernel} we have
  \begin{center}
    \begin{tikzcd}
      A \arrow[r, "\varphi"] & B \arrow[dr, "\coker \varphi"'] \arrow[rr, "\gamma \circ \coker \varphi = 0"] &   & Z \\
        &   & C \arrow[ur, bend left, "\gamma"'] \arrow[ur, dotted, bend right, "\exists!"'] &
    \end{tikzcd}
  \end{center}
  Therefore since \(\gamma \circ \coker \varphi = 0 = 0 \circ \coker \varphi\) uniqueness forces \(\gamma = 0\) and thus by \autoref{lemma:mono-epi-additive} \(\varphi\) is an \hyperref[def:mono-epi-morphism]{epimorphism}.
\end{proof}

\begin{lemma} \label{lemma:partial-sequence-additive}
  Let \(\varphi: A \to B\) be a morphism in an \hyperref[def:additive-category]{additive category}. Then \(\varphi\) is a \hyperref[def:mono-epi-morphism]{monomorphism} if and only if \(0 \to A\) is its \hyperref[def:categorical-kernel]{kernel} and \(\varphi\) is an \hyperref[def:mono-epi-morphism]{epimorphism} if and only if \(B \to 0\) is its \hyperref[def:categorical-cokernel]{cokernel}.
\end{lemma}
\begin{proof}
  We will do the \hyperref[def:categorical-kernel]{kernel} one.
  \begin{itemize}
    \item[\(\implies\)] Suppose \(\varphi\) is a \hyperref[def:mono-epi-morphism]{monomorphism}.

          If \(\xi: Z \to A\) is a morphism such that \(\varphi \circ \xi = 0\) then by \autoref{lemma:mono-epi-additive} \(\xi = 0\).
          In particular \(\xi\) factors uniquely through \(0 \to A\), i.e: \(0 \to A\) is the \hyperref[def:categorical-kernel]{kernel} of \(\varphi\).
    \item[\(\impliedby\)] Suppose now that \(0 \to A\) is the \hyperref[def:categorical-kernel]{kernel} of \(\varphi\).

          Consider a morphism \(\xi: Z \to A\) such that \(\varphi \circ \xi = 0\).
          By the \hyperref[def:categorical-kernel]{kernel} properties we have the following diagram
          \begin{center}
            \begin{tikzcd}
              Z \arrow[rr, "\xi"'] \arrow[rrr, bend left, "0"] \arrow[dr, dotted, "\exists!"'] & & A \arrow[r, "\varphi"'] & B \\
              & 0 \arrow[ur]
            \end{tikzcd}
          \end{center}
          Therefore by uniqueness \(\xi = 0\), making \(\varphi\) a \hyperref[def:mono-epi-morphism]{monomorphism} by \autoref{lemma:mono-epi-additive}.
  \end{itemize}
\end{proof}

\subsection{Additive functors}
\begin{definition}[Additive functor] \label{def:additive-functor}
  Let \(\category{A}, \category{B}\) be two \hyperref[def:additive-category]{additive categories}. A \hyperref[def:functor]{functor} \(F: \category{A} \to \category{B}\) is \textbf{additive} if it preserves the abelian group structure of all Hom-sets.
\end{definition}

\begin{lemma}\label{lemma:additive-functor-preserves-sum}
  Let \(F: \category{A} \to \category{B}\) be a \hyperref[def:functor]{functor} between two \hyperref[def:additive-category]{additive categories} \(\category{A}\) and \(\category{B}\).
  Then \(F\) transforms \hyperref[def:direct-sum]{direct sums} to \hyperref[def:direct-sum]{direct sums}.
\end{lemma}
\begin{proof}
  Let \(X, Y \in \category{A}\) be two objects and \(X \oplus Y\) be its \hyperref[def:direct-sum]{direct sum}.
  Applying \(F\) to the diagram in \autoref{rem:morphisms-in-direct-sum} defining \(A \oplus B\) we get
  \begin{center}
    \begin{tikzcd}
      F(X) \arrow[rr, "\id_{F(X)}"] \arrow[dr, "Fi"] & & F(X) \\
      & F(X \oplus Y) \arrow[ur, "Fp"] \arrow[dr, "Fq"] & \\
      F(Y) \arrow[ur, "Fj"] \arrow[rr, "\id_{F(Y)}"] & & F(Y)
    \end{tikzcd}
  \end{center}
  Note that
  \begin{align*}
    Fp \circ Fj &= F(p \circ j) = F(0) = 0 \\
    Fq \circ Fi &= F(q \circ i) = F(0) = 0
  \end{align*}
  Furthermore since \(F\) is \hyperref[def:additive-functor]{additive}
  \begin{equation*}
    Fp \circ Fi + Fj \circ Fq = F(p \circ i + j \circ q) = F\id_{X \oplus Y} = \id_{F(X \oplus Y)}
  \end{equation*}
  Therefore \(F(X \oplus Y)\) is also a \hyperref[def:direct-sum]{direct sum}.
\end{proof}


\begin{lemma}\label{lemma:additive-functor-characterization}
  Let \(F: \category{A} \to \category{B}\) be a \hyperref[def:functor]{functor} between two \hyperref[def:additive-category]{additive categories} \(\category{A}\) and \(\category{B}\). The following are equivalent:
  \begin{enumerate}
    \item \(F\) is \hyperref[def:additive-functor]{additive}.
    \item For all \(A, B \in \category{A}\) we have an \hyperref[def:categorical-isomorphism]{isomorphism}
          \begin{equation*}
            F(A) \oplus F(B) \isoarrow F(A \oplus B)
          \end{equation*}
    \item For all \(A, B \in \category{A}\) we have an \hyperref[def:categorical-isomorphism]{isomorphism}
          \begin{equation*}
             F(A \oplus B) \isoarrow F(A) \oplus F(B)
          \end{equation*}
  \end{enumerate}
\end{lemma}
\begin{proof}
  Since by \autoref{lemma:additive-functor-preserves-sum} \(F\) commutes with \hyperref[def:direct-sum]{direct sums} 1 implies 2 and 3.
  Furthermore since the composition
  \begin{center}
    \begin{tikzcd}
      F(A) \oplus F(B) \arrow[r] & F(A \oplus B) \arrow[r] & F(A) \oplus F(B)
    \end{tikzcd}
  \end{center}
  is the identity map 2 and 3 are equivalent.

  To finish the proof, assume 2 and 3 hold and consider two maps \(f, g: A \to B\).
  Note that \(f + g\) is equal to the composition
  \begin{center}
    \begin{tikzcd}
      A \arrow[r] & A \oplus A \arrow[r, "f \oplus g"] & B \oplus B \arrow[r] & B
    \end{tikzcd}
  \end{center}
  Consider the diagram
  \begin{center}
    \begin{tikzcd}
      F(A) \arrow[dr] \arrow[r] & F(A \oplus A) \arrow[r, "F(f \oplus g)"] & F(B \oplus B) \arrow[r] \arrow[d, "\sim"] & F(B) \\
      &  F(A) \oplus F(A) \arrow[r, "Ff \oplus Fg"] \arrow[u, "\sim"] &  F(B) \oplus F(B) \arrow[ur]
    \end{tikzcd}
  \end{center}
  Checking for each of the four induced maps \(F(A) \to F(B)\) since \(F\) is a \hyperref[def:functor]{functor} we have that the middle square commutes.
  Moreover since the vertical maps are \hyperref[def:categorical-isomorphism]{isomorphisms} the triangles commute.
  Therefore we get
  \begin{equation*}
    Ff \oplus Fg = F(f \oplus g)
  \end{equation*}
\end{proof}

\end{document}
